% WHAT A MARK SHOWS OF ITSELF.
%
% \show and \meaning give one of TeX's five marks its name, a colon, and
% then the TEXT it holds -- the only primitives besides a macro that show
% more than a name:
%
%   \show\botmark        > \botmark=\botmark:
%                        om##ega.
%   \meaning\botmark     \botmark:om##ega
%
% \show breaks the line after the colon, the way it breaks a macro's
% between what it is and what it says; \meaning does not, for the same
% reason it does not break a macro's.
%
% The text is written the way a token list is written anywhere, so a
% parameter character in it is doubled.
%
% The marks CLASSES eTeX adds do NOT do this: `\show\botmarks' is
% `> \botmarks=\botmarks.' and no more.
%
% A TRACE, though, gives the name alone: \tracingcommands writes
% `{\botmark}' and \tracingrestores writes `{restoring \botmark=\botmark}'.
% Only \show and \meaning ask what a command HOLDS; a trace asks what it
% IS.
%
% HSTeX gave the name alone everywhere, and then, once \show was put right,
% gave the text everywhere. trip line 320 writes \show\botmark.
%
% The text a mark expands into is read back from a frame the reference
% calls `<mark>', not `<inserted text>'.
%
% And a parameter character with nothing to do is TRACED before it is
% refused: `{macro parameter character #}' and then the fault. The trace
% says what the main loop was given, whatever it then makes of it.
\catcode`\{=1 \catcode`\}=2 \catcode`\#=6
\tracingonline=1
\output={\global\setbox9=\box255}
\vsize=10pt \hsize=100pt
\message{[A every mark is empty to begin with]}
\show\botmark
\show\topmark
\show\splitbotmark
\mark{alpha}\hrule height 5pt \mark{om#ega}\hrule height 20pt
\penalty-10000
\message{[B a page has been made, so the marks hold what it left]}
\show\topmark
\show\firstmark
\show\botmark
\message{[C the same three, asked for as a meaning]}
\message{<\meaning\topmark><\meaning\firstmark><\meaning\botmark>}
\message{[D a \let to one shows the mark it names]}
\let\bm=\botmark \show\bm
\message{[E a trace names a mark without its text]}
\tracingcommands=2
\botmark
\let\bm=\botmark \bm
\tracingcommands=0
\message{[F so does a restore]}
\tracingrestores=1
\begingroup \let\botmark=\relax \endgroup
\tracingrestores=0
\message{[done]}
\end
